% Source printed pp. 181-185 / PDF pp. 189-193.
\section{Deligne--Mumford 叠的局部结构}\label{src:5.3}
\begin{epigraph}
要是我有了那些定理就好了! 那么我就能轻易找出证明.
\epigraphauthor{相传出自 Bernhard Riemann}
\end{epigraph}
{\emergencystretch=3em
我们将说明, Deligne--Mumford 叠 $\mathcal X$ 在点 $x$ 附近, 在平展局部是仿射概形对稳定子群概形的商叠 $[\Spec A/G_x]$. 从概念上说, 这表明, 正如概形 (分别为代数空间) 由仿射概形在 Zariski 拓扑 (分别为平展拓扑) 中粘合而成, Deligne--Mumford 叠由商叠 $[\Spec A/G]$ 在平展拓扑中粘合而成.\footnote{当然, Deligne--Mumford 叠在平展局部也是概形, 但定理~\bookref{5.3.1} 所构造的平展邻域 $([\Spec A/G_x],w)\to(\mathcal X,x)$ 保持 $w$ 处的稳定子群.} 这一点的实际用途是, 它使我们能够将 Deligne--Mumford 叠的许多性质归结到商叠 $[\Spec A/G]$. 我们将利用局部结构定理构造粗模空间 (定理~\bookref{5.4.6}).\par}


\subsection{局部结构定理}\label{src:5.3.1-subsection}
Deligne--Mumford 叠 $\mathcal X$ 的点 $x\in|\mathcal X|$ 的\newterm{几何稳定子}{geometric stabilizer} 是如下定义的抽象群: 取任意像为 $x$ 的几何点 $\Spec k\to\mathcal X$, 并取其稳定子.
\begin{theorem}[Deligne--Mumford 叠的局部结构定理]\label{src:5.3.1}
设 $\mathcal X$ 是拟分离的 Deligne--Mumford 叠, $x\in|\mathcal X|$ 是有限型点, 几何稳定子为 $G_x$. 存在平展态射
\[
f:([\Spec A/G_x],w)\longrightarrow(\mathcal X,x),
\]
其中 $w\in[\Spec A/G_x]$, 使得 $f$ 在 $w$ 处诱导几何稳定子群的同构. 此外, 若 $\mathcal X$ 的对角线分离 (分别为仿射), 则可使 $f$ 可表出且分离 (分别为仿射).
\end{theorem}
\begin{proof}
先处理 $\mathcal X$ 的对角线仿射的情形. 选取代表 $x$ 的几何点 $\Spec k\to\mathcal X$, 设 $d$ 为 $G_x$ 的基数. 将 $G_x$ 看作 $\Spec k$ 上的群概形, 设 $BG_x\to\mathcal X$ 为所诱导的映射. 设 $(U,u)\to(\mathcal X,x)$ 是来自仿射概形的平展可表出态射. 由于 $\mathcal X$ 的对角线仿射, $U\to\mathcal X$ 仿射. 定义拟仿射概形
\[
\mathrm{SEC}_d:=\underbrace{U\times_{\mathcal X}\cdots\times_{\mathcal X}U}_{d\text{ 次}}\setminus\Delta,
\]
其中 $\Delta$ 是所有两两对角线的并. 来自概形的映射 $S\to\mathrm{SEC}_d$ 由一个态射 $S\to\mathcal X$ 以及 $U_S:=U\times_{\mathcal X}S\to S$ 的 $d$ 个不交截面 $s_1,\ldots,s_d$ (即当 $i\neq j$ 时, $s_i$ 与 $s_j$ 的交为空) 分类. 置换截面给出 $S_d$ 在 $\mathrm{SEC}_d$ 上的作用, 定义商叠
\[
\mathrm{ET}_d:=[\mathrm{SEC}_d/S_d].
\]
由主 $S_d$-丛与次数为 $d$ 的有限平展覆盖之间的对应 (习题~\bookref{B.1.55}), $\mathrm{ET}_d$ 在概形 $S$ 上的对象对应于图
\[
\begin{tikzcd}
Z\arrow[r,hook]\arrow[dr]&U_S\arrow[r]\arrow[d]&U\arrow[d]\\
&S\arrow[r]\arrow[ur,phantom,"\square" description]&\mathcal X,
\end{tikzcd}
\]
其中 $Z\hookrightarrow U_S$ 是闭子概形, $Z\to S$ 是次数为 $d$ 的有限平展态射. 有所诱导的平展态射 $\mathrm{ET}_d\to\mathcal X$ 以及交换图
\[
\begin{tikzcd}[column sep=large]
\mathrm{SEC}_d\arrow[r]\arrow[d,hook]&\mathrm{ET}_d\arrow[dr]&\\
U\times_{\mathcal X}\cdots\times_{\mathcal X}U\arrow[r]&U\arrow[r]&\mathcal X.
\end{tikzcd}
\]
为构造 $x$ 的一个原像 $w\in\mathrm{ET}_d(k)$, 注意纤维积 $BG_x\times_{\mathcal X}U$ 是有限个 $\Spec k$ 的不交并, 可以选取其中一个. 这给出交换图
\[
\begin{tikzcd}
\Spec k\arrow[r,hook]\arrow[dr]&\coprod\Spec k\arrow[r]\arrow[d]&U\arrow[d]\\
&BG_x\arrow[r]\arrow[ur,phantom,"\square" description]&\mathcal X,
\end{tikzcd}
\]
其中 $\Spec k\to BG_x$ 是次数为 $d$ 的有限平展态射, 这定义了 $BG_x\to\mathcal X$ 的提升 $BG_x\to\mathrm{ET}_d$. 复合 $w:\Spec k\to BG_x\to\mathrm{ET}_d$ 是 $x$ 的原像. 我们断言映射 $\mathrm{ET}_d\to\mathcal X$ 在 $w$ 处诱导几何稳定子的同构. 为说明这一点, 只需假设 $\mathcal X=BG_x$, $U=\Spec k$. 此时有同构 $U\times_{\mathcal X}\cdots\times_{\mathcal X}U\cong G_x^{d-1}$, 还可将其与对角作用的商 $G_x^d/G_x$ 等同. 于是 $\mathrm{SEC}_d\subseteq U\times_{\mathcal X}\cdots\times_{\mathcal X}U$ 与商 $(G_x^d\setminus\Delta)/G_x$ 等同. 置换作用 $S_d$ 是传递的, 其稳定子同构于 $G_x$. 事实上, $\mathrm{SEC}_d$ 与正则表示 $G\subseteq S_d$ 的商 $S_d/G$ 是 $S_d$-等变同构的, 因而 $\mathrm{ET}_d\cong BG_x$.

由于 $\mathrm{ET}_d$ 分离且 $\mathcal X$ 的对角线分离, 相对惯性叠 $I_{\mathrm{ET}_d/\mathcal X}\to\mathrm{ET}_d$ 有限. 由于 $\mathrm{ET}_d\to\mathcal X$ 在 $w$ 处诱导稳定子的同构, $I_{\mathrm{ET}_d/\mathcal X}\to\mathrm{ET}_d$ 在 $w$ 附近的某个开子叠 $[W/S_d]\subseteq\mathrm{ET}_d$ 上是同构, 其中 $W\subseteq\mathrm{SEC}_d$ 是拟仿射概形. 因此 $[W/S_d]\to\mathcal X$ 是分离, 平展且可表出的态射, 在 $w$ 处诱导稳定子群的同构. 改为对 $G_x\subseteq S_d$ 取商后, $[W/G_x]\to\mathcal X$ 同样是分离, 平展且可表出的态射, 在 $w$ 处诱导稳定子群的同构. 设 $W'\subseteq W$ 是包含 $w$ 的仿射开子概形, 可将 $W$ 换成 $G_x$-不变的仿射开子概形 $\bigcap_{g\in G_x}g\cdot W'$. 还需说明 $[W/G_x]\to\mathcal X$ 仿射. 由于 $\mathcal X$ 的对角线仿射, 来自仿射概形 $W$ 的态射 $W\to\mathcal X$ 仿射. 纤维积
\[
\begin{tikzcd}
{[W/G_x]}\times_{\mathcal X}W\arrow[r]\arrow[d]&W\arrow[d]\\
{[W/G_x]}\arrow[r]\arrow[ur,phantom,"\square" description]&\mathcal X
\end{tikzcd}
\]
在 $[W/G_x]$ 上仿射, 因而同构于商叠 $[\Spec B/G_x]$. 另一方面, 由于 $[W/G_x]\to\mathcal X$ 可表出, 商叠 $[\Spec B/G_x]$ 是代数空间, 且 $G_x$ 在 $\Spec B$ 上的作用自由. 由推论~\bookref{5.2.13}, $[\Spec B/G_x]$ 同构于仿射概形 $\Spec B^{G_x}$. 由平展下降, $[W/G_x]\to\mathcal X$ 仿射.

若 $\mathcal X$ 的对角线分离, 则上述论证给出具有 $G_x$ 作用的分离概形 $\mathrm{SEC}_d$, 以及分离, 平展且可表出的态射 $([\mathrm{SEC}_d/G_x],w)\to(\mathcal X,x)$, 在 $w$ 处诱导同构. 将论证再次用于分离的 Deligne--Mumford 叠 $[\mathrm{SEC}_d/G_x]$, 即得所需的局部商呈示. 若 $\mathcal X$ 仅拟分离, 则论证给出态射 $([\mathrm{SEC}_d/G_x],w)\to(\mathcal X,x)$ (不一定可表出), 在 $w$ 处诱导同构, 其中 $\mathrm{SEC}_d$ 是拟分离代数空间. 由于 $\mathrm{SEC}_d$ 的对角线分离 (因对角线是单态射), $[\mathrm{SEC}_d/S_d]$ 是具有分离对角线的 Deligne--Mumford 叠, 可以应用对角线分离情形的结论, 得到由仿射概形的商给出的局部商呈示. 另见~\cite[定理 $6.2$]{LMB00}.
\end{proof}
\begin{exercise}\label{src:5.3.2}
设 $\mathcal X$ 是 Deligne--Mumford 叠. 证明 $\mathcal X$ 同构于商叠 $[U/G]$, 其中 $U$ 是仿射概形 (分别为概形, 代数空间), $G$ 是有限抽象群, 当且仅当存在来自仿射概形 (分别为概形, 代数空间) 的有限平展态射 $V\to\mathcal X$.

\textit{提示:} 若 $V\to\mathcal X$ 是次数为 $d$ 的有限平展覆盖, 考虑如习题~\bookref{B.1.55} 中所述的相应主 $S_d$-torsor
\[
\underbrace{V\times_{\mathcal X}\cdots\times_{\mathcal X}V}_{d\text{ 次}}\setminus\Delta\longrightarrow\mathcal X.
\]
\end{exercise}
\begin{proposition}\label{src:5.3.3}
若 $R\rightrightarrows U$ 是仿射概形的有限平展等价关系, 则代数空间商 $U/R$ 是仿射概形.
\end{proposition}
\begin{proof}
由习题~\bookref{5.3.2}, 代数空间 $U/R$ 同构于 $V/G$, 其中有限群 $G$ 在仿射概形 $V=\Spec B$ 上自由作用. 定理~\bookref{5.2.3} 表明 $V/G\to\Spec B^G$ 对到代数空间的映射具有泛性质, 因而是同构. 也可由习题~\bookref{5.2.16} 推出: 若 $U=\Spec A$, 则 $U/R\to\Spec A^R$ 对到代数空间的映射具有泛性质, 因而是同构.
\end{proof}

\subsection{将有理点提升到呈示}\label{src:5.3.2-subsection}
使用与局部结构定理 (\bookref{5.3.1}) 的证明类似的技巧, 可证明存在能提升给定域值点的呈示.
\begin{proposition}\label{src:5.3.4}
若 $\mathcal X$ 是具有分离对角线的代数叠 (分别为具有分离且拟紧对角线的 Deligne--Mumford 叠), $x\in\mathcal X(k)$ 是域值点, 则存在来自仿射概形的光滑 (分别为平展) 态射 $U\to\mathcal X$, 以及位于 $x$ 上方的点 $u\in U(k)$.
\end{proposition}
\begin{proof}
设 $U\to\mathcal X$ 是来自仿射概形的光滑态射, 其像包含 $x$. 考虑纤维积
\[
\begin{tikzcd}
U_x\arrow[r]\arrow[d]&U\arrow[d]\\
\Spec k\arrow[r,"x"']\arrow[ur,phantom,"\square" description]&\mathcal X.
\end{tikzcd}
\]
由于 $\mathcal X$ 的对角线分离, 态射 $U\to\mathcal X$ 分离. 由于 $U_x$ 在 $\Spec k$ 上有限型且拟分离, 它是 Noether 代数空间. 选取闭点 $u\in U_x$, 由剩余 gerbe 的存在性 (\bookref{4.5.16}), 存在来自剩余域的闭浸入 $\Spec\kappa(u)\to U_x$. 域扩张 $k\to\kappa(u)$ 有限且可分, 设其次数为 $d$.

沿用定理~\bookref{5.3.1} 证明中的记号, 考虑参数化 $U\to\mathcal X$ 的 $d$ 个不交截面的开子空间 $\mathrm{SEC}_d\subseteq U\times_{\mathcal X}\cdots\times_{\mathcal X}U$, 它是分离代数空间. 定义分离的 Deligne--Mumford 叠 $\mathrm{ET}_d=[\mathrm{SEC}_d/S_d]$, 并考虑所诱导的态射 $\mathrm{ET}_d\to\mathcal X$. 由于 $\Spec\kappa(u)\to\Spec k$ 是次数为 $d$ 的有限平展态射, 闭浸入 $\Spec\kappa(u)\hookrightarrow U_x$ 定义 $\mathrm{ET}_d$ 的一个 $k$-点 $v$. 这给出交换图
\[
\begin{tikzcd}
&&\Spec k\arrow[dl,"v"']\arrow[d,"x"]\\
\mathrm{SEC}_d\arrow[r]&\mathrm{ET}_d\arrow[r]&\mathcal X.
\end{tikzcd}
\]
点 $v\in\mathrm{ET}_d(k)$ 不一定提升到 $\mathrm{SEC}_d$, 但可以用下面的技巧选取 $\mathrm{ET}_d$ 的另一种商呈示, 使 $v$ 可以提升.

具体而言, 选取忠实表示 $S_d\subseteq\GL_n$, 写成 $\mathrm{ET}_d\cong[V/\GL_n]$, 其中 $V=\mathrm{SEC}_d\times^{S_d}\GL_n$ 是分离代数空间. 于是 $v:\Spec k\to[V/\GL_n]$ 对应于主 $\GL_n$-丛 $P\to\Spec k$ 以及 $\GL_n$-等变映射 $P\to V$. 由于主 $\GL_n$-丛与向量丛一一对应 (习题~\bookref{B.1.58}), $P$ 是平凡主 $\GL_n$-丛, 存在截面 $\Spec k\to P$. 复合 $V\to[V/\GL_n]\to\mathcal X$ 光滑, 复合 $\Spec k\to P\to V$ 是 $x$ 的提升.

由于 $V$ 是分离代数空间, 问题归结为证明第二个陈述, 即具有分离且拟紧对角线的 Deligne--Mumford 叠 $\mathcal X$ 的 $k$-点可提升到来自概形的平展邻域. 将用到 $\mathcal X$ 的对角线拟仿射这一事实, 它将在后面的推论~\bookref{5.5.8} 中独立证明 (证明只用到上面的命题~\bookref{5.3.3} 以及~\S\bookref{5.1} 中的拟凝聚层理论). 重复上述论证, 选取来自仿射概形的平展映射 $U\to\mathcal X$, 其像包含 $x$. 则 $U\to\mathcal X$ 的 $d$ 个不交截面的空间 $\mathrm{SEC}_d$ 是拟仿射概形, 且带有 $S_d$ 的自由作用. 商 $\mathrm{ET}_d=\mathrm{SEC}_d/S_d$ 也是拟仿射的 (习题~\bookref{5.2.14}). 所诱导的映射 $\mathrm{ET}_d\to\mathcal X$ 平展, 且按构造 $k$-点 $x$ 提升为 $\mathrm{ET}_d$ 的 $k$-点. 另见~\cite[定理 $6.3$]{LMB00}.
\end{proof}
还有局部结构定理的另一个版本, 在粗模空间上以平展局部的方式表述 (见推论~\bookref{5.4.13}).
